\chapter*{Introduction}

Historiquement, la logique et l'algèbre se sont développées comme deux
disciplines distinctes avec peu d'interactions. Cependant, le développement de
la théorie des modèles permet de créer un lien entre ces deux domaines des mathématiques.

L'objectif principal de ce mémoire est d'illustrer ce lien en étudiant en
particulier les corps algébriquement clos. L'enjeu est de montrer comment des
notions purement logiques comme théorème de compacité ou l'élimination des
quantificateurs permettent de résoudre de manière efficace et élégante des problèmes de
nature algébrique. En particulier, on va utiliser la théorie des modèles pour démontrer deux résultats
classiques : d'Ax-Grothendieck, qui affirme que toute application polynomiale
injective de $\bC^n$ dans lui-même est surjective, et le célèbre Nullstellensatz
de Hilbert.

Avant d'entrer dans le vif du sujet, nous voulons remercier chaleureusement Adam
Parusiński, pour avoir accepté de nous encadrer, de nous avoir fait découvrir ce
sujet original et stimulant, ainsi que pour ses conseils et ses relectures tout
le long du projet.

Afin d'atteindre notre objectif, nous avons organisé ce mémoire en trois
chapitres:

\begin{itemize}
    \item Le chapitre 1 donne les rudiments de la théorie des modèles, suivant de près
    le plan de \cite{Marker_2002}, dans lequel on pourra trouver plus de détails. Nous
    introduisons le formalisme et les résultats nécessaire pour la suite, en particulier 
    le théorème de compacité et le test de Vaught. On introduit aussi les la
    théorie des corps algébriquement clos (ACF).
    \item est consacré à l'étude des corps algébriquement clos sous l'angle de
    la complétude. 
    Grâce à laquelle nous y établissons le principe de transfert de Lefschetz, dont nous illustrons la puissance en démontrant le théorème géométrique d'Ax-Grothendieck.
    \item Le chapitre 3 aborde l'élimination des quantificateurs qui nous permet de lier la logique à la géométrie via les ensembles constructibles, aboutissant à une preuve élégante du Nullstellensatz.
\end{itemize}

\addcontentsline{toc}{chapter}{Introduction}
